\subsection{Pregunta 5: impacto del ancho de banda en las métricas OSPF}
\autor{Sebastián}

\textbf{Escenario:} en el enlace Deadlock--Raze (2.0.0.0) se configura \cmd{bandwidth 150} y se envía un paquete de Annihilation a Paint Shells.

El costo OSPF de una interfaz es
\[
  \text{Costo} = \frac{10^{8}}{\text{Bandwidth (bps)}}
\]
con parte entera y ancho de banda de referencia de 100~Mbps, el valor por defecto de Cisco.

\subsubsection*{Cálculo de costos}
\begin{itemize}[nosep]
  \item Enlace directo con \cmd{bandwidth 150}: $\left\lfloor \dfrac{10^{8}}{150\,000} \right\rfloor = 666$, más el costo de la interfaz LAN de Raze (\pend{valor}).
  \item Trayecto alternativo Deadlock $\rightarrow$ Cypher $\rightarrow$ Killjoy $\rightarrow$ Raze: \pend{suma de los costos de cada enlace, según \cmd{show ip ospf interface}}.
\end{itemize}

\evid{Q5_ancho_banda/show_ip_route_deadlock.png}{Tabla de rutas de Deadlock con el ancho de banda modificado.}
\evid{Q5_ancho_banda/tracert_annihilation_paint.png}{Traza de Annihilation a Paint Shells.}

\subsubsection*{Conclusión}
\pend{¿el tráfico sigue por 2.0.0.0 o conmuta al trayecto alternativo? Corroborar con la tabla de rutas}

\textit{Recordatorio: restaurar el ancho de banda original después de la prueba.}
